% concatenated from CCV.tex
\documentclass[a4paper]{amsart}
\usepackage{url}
\usepackage{amsmath}
\usepackage[makeroom]{cancel}

%\usepackage{refcheck}
%\usepackage[notref,notcite]{showkeys}

\usepackage{graphicx}
\usepackage[all]{xy}
\usepackage[mathscr]{eucal}
\usepackage{amsmath,amssymb,amsfonts}

%\usepackage{mathrsfs,latexsym,amsthm,enumerate}
%\usepackage{amscd}

\newtheorem{theorem}{Theorem}[section]
\newtheorem{lemma}[theorem]{Lemma}
\newtheorem{corollary}[theorem]{Corollary}
\newtheorem{example}[theorem]{Example}
\newtheorem{examples}[theorem]{Examples}
\newtheorem{proposition}[theorem]{Proposition}
\newtheorem{remark}[theorem]{Remark}
\newtheorem{conjecture}[theorem]{Conjecture}
\newtheorem{question}[theorem]{Question}
\newtheorem{terminology}[theorem]{Terminology}

\usepackage{stackengine}
\usepackage{xcolor}
\usepackage[all]{xy}


\title{Finitary Cartesian closed varieties and semigroup actions}

\author{Mark V. Lawson}
\address{Mark V. Lawson, Department of Mathematics,
Maxwell Institute for Mathematical Sciences,
Heriot-Watt University,
Riccarton,
Edinburgh EH14 4AS,
UNITED KINGDOM}
\email{m.v.lawson@hw.ac.uk}



\begin{document}
\dedicatory{This paper is dedicated to the memory of Philip Scott, friend and colleague.}


%%%%%%%%%%%%%%%%%%%%%%%%%%%%%%%%%%%%%%%%%%%%%%%%%%%%%%%%%%%%%%%%%%%%%%%%%%%%%%%%%%%%%%%%%%%%%%
\begin{abstract}
We build on some ideas of Richard Garner.
Let $M$ be a monoid and $B$ a Boolean algebra.
A `matched pair' $[B|M]$ consists of $B$ and $M$ and some mutual interactions.
Garner showed that every such matched pair determines (what we shall call)
a Boolean left restriction monoid $S = S[B|M]$.
In this paper, we show that the data of a $[B|M]$-set (defined later)
may be encoded by means of a certain kind of action by $S$.
This means that the category $[B|M]$-{\bf sets} is equivalent to
a category of {\bf $S$-actions}.
We deduce, as a result of Garner's work, that every non-degenerate finitary Cartesian closed variety 
is equivalent to a special category of $S$-actions where $S$ is a Boolean left restriction monoid.

\end{abstract}
\maketitle


%%%%%%%%%%%%%%%%%%%%%%%%%%%%%%%%%%%%%%%%%%%%%%%%%%%%%%%%%%%%%%%%%%%%%%%%%%%%%%%%%%%%%%%%%%%%%%%%%%
\section{Introduction}

This paper builds on some excellent work of Richard Garner  \cite{G2024, G2025}.
Before we explain what we actually do, we describe our point of view; this will serve to orientate the reader.

Semigroup theory is, of course, defined as the study of associative binary operations.
In practice, however, it can often be regarded as the {\em algebraic} study of functions, to be contrasted with
category theory and the theory of computable functions.
The subtitle of \cite{JS} describes what we have in mind.
Just as with category theory, the word `function' should be interpreted broadly.
Thus, inverse semigroup theory is the abstract theory of {\em partial} bijections.
Similarly, the left restriction monoids, which play a starring role in this paper,
should be viewed as the abstract theory of partial functions along with their domains of definition;
to be precise, you should think in terms of partial functions as being composed from left-to-right.
In addition, we shall need a class of left actions of tubesuch monoids.
As we shall see, you can regard such actions as being on presheaves of sets. 

In Section~4, we prove that a category consisting of a single Boolean left restriction monoid $S$
acting on supported sets
is equivalent to the category 
$[\mathsf{Proj}(S)|\mathsf{Tot}(S)]$-{\bf set}
(we shall explain all terminology and notation later).
However, Garner \cite{G2024, G2025} proved that 
every non-degnerate finitary Cartesian closed variety is isomorphic to one of the form {\bf $[B|M]$-sets}
where $B$ is a Boolean algebra and $M$ is a monoid.
It will therefore follow from our work that every Cartesian closed variety 
is equivalent to a category consisting of a single Boolean left restriction monoid $S$
acting on supported sets.

Why should you care? The groups of units of what are termed `\'etale Boolean left restriction monoids'
are topological full groups of the `ample groupoids' within the theory of $C^{\ast}$-algebras.
Such groups can be seen as generalizations of the Thompson groups, which were originally constructed from the $\lambda$-calculus
(whose categorical incarnation is the the theory of Cartesian closed categories).

The reference \cite{L1988} will not actually be needed to read this paper,
but I have included it since Phil Scott would have been amused to see a connection between
Cartesian closed categories and semigroup theory.

The paper proceeds as follows.
In Section~2,
we describe a class of left restriction monoids that we call `factorizable'
in terms of `matched pairs' $[E|M]$ where here $E$ is a meet semilattice with a top element
and $M$ is a monoid. 
It will turn out that Boolean left restriction monoids are always factorizable.
In Section~3,
we introduce a class of actions of left restriction monoids that we term `supported';
this generalizes the \'etale actions of inverse monoids.
We also describe a category obtained from $[E|M]$ that we call the category of $[E|M]$-{\bf sets}.
We then fix a left restriction monoid $S$ and prove that 
a class of supported actions of $S$, that we also term `factorizable', is equivalent to the category of $[E|M]$-{\bf sets}.
In Section~4, we show how the above theory works out when we specialize to  Boolean
left restriction monoids; it is here that we meet Garner's work again.
In Section~5, we describe the results of Section 4 in the case of those Boolean left restriction monoids which are \'etale.
Thus we proceed from the general to the specific.\\

\noindent
{\bf Terminology and notation}\\
The word `semilattice' will mean a meet semilattice with a top element denoted by 1.
The meet operation will be denoted by concatenation.
We assume the reader is familar with Boolean algebras.
Furthermore,  we shall assume that $0 \neq 1$ (so, our Boolean algebras are `non-degenerate').
The complement of $e$ will be denoted by $\bar{e}$,
the meet of $e$ and $f$ by $ef$
and the join of $e$ and $f$ by $e + f$.
I shall use the word `homomorphism' to mean any structure-preserving function.
Thus homomorphisms of monoids will be required to preserve the identity
and homomorphisms of semigroups-with-zeros will be required to preserve the zero.\\


\noindent
{\bf Acknowledgements }The author would like to thank the Isaac Newton Institute for Mathematical Science, Cambridge,
for support and hospitality during the programme {\em Topological groupoids and $C^{\ast}$-algebras},
where work on this paper was undertaken.
This work was supported by EPSRC grant EP/V521929/1.
The author would particularly like to thank
Ying-Fen Lin (Queen's, Belfast) for the invitation
as well as the other organizers
Astrid an Huef (Victoria, Wellington),
Nadia Larsen (Oslo),
Xin Lin (Glasgow),
and
Stuart White (Oxford).\\





%%%%%%%%%%%%%%%%%%%%%%%%%%%%%%%%%%%%%%%%%%%%%%%%%%%%%%%%%%%%%%%%%%%%%%%%%%%%%%%%%%%%%%%%%%%%%%%%%%%%%%%%%%%%%%%%%%%%%%%%%%%%%%%%%%%%%%%%%
\section{The structure of factorizable  left restriction monoids}

The goal of this section is to describe the structure of what we call `factorizable left restriction monoids'
in terms of semilattices and monoids using ideas to be found in \cite{G2025}.
It will transpire that the Boolean left restriction monoids, which we are really interested in,
are automatically factorizable in our sense.

Define a monoid $S$ to be a {\em left restriction monoid} if
it is equipped with a unary operation $s \mapsto s^{+}$, called {\em plus}, satisfying the following axioms:
\begin{itemize}
\item[{\rm (LR1).}] $(s^{+})^{+} = s^{+}$. 
\item[{\rm (LR2).}] $(s^{+}t^{+})^{+} = s^{+}t^{+}$.
\item[{\rm (LR3).}] $s^{+}t^{+} = t^{+}s^{+}$.
\item[{\rm (LR4).}] $s^{+}s = s$.
\item[{\rm (LR5).}] $(st)^{+} = (st^{+})^{+}$.
\item[{\rm (LR6).}] $st^{+} = (st)^{+}s$.
\end{itemize}
See my paper \cite{L2025} for the origins of these axioms (or, rather, their left-right duals)
and the basic properties that follow.

Denote by $\mathsf{Proj}(S)$ those elements $a$ such that $a^{+} = a$, called {\em projections}.
This set is a meet semilattice with top element (in other words, it is a monoid under meet).
Those elements $a \in S$ such that $a^{+} = 1$ are called  {\em total elements};
the set of total elements of $S$ is denoted by $\mathsf{Tot}(S)$.
This is a monoid by axiom (LR5).

Define a partial order on a left restriction monoid by $s \leq t$ iff $s = s^{+}t$.
We call this the {\em natural partial order} and it will be the only order we shall consider
on a left restriction monoid.
It can be shown that $s \leq t$ iff $s = et$ where $e$ is a projection.
If $s \leq t$ then $s^{+} \leq t^{+}$;
we shall use this result many times.


Let $S$ and $T$ be left restriction monoids.
A {\em homomorphism} $\theta \colon S \rightarrow T$ of left restriction monoids
is a monoid homomorphism from $S$ to $T$ such that $\theta (a^{+}) = \theta (a)^{+}$.
Such homomorphisms map projections to projections.

We shall now define a structure form any left restriction monoid.

%%%%%%%%%%%%%%%%%%%%%%%%%%%%%%%%%%%%%%%%%%%%%%%%%%%%%%%%%%%%%%%%%%%%%%%%%%%%%%%%%%%%%%%%%%%%%%%%%%%%%%%%%%%%%%%%%%%%%%%%%%%
Put $M = \mathsf{Tot}(S)$ and $E = \mathsf{Proj}(S)$.
We shall denote the product in $M$ by concatenation as well as the meet in $E$.
Define a function $M \times E \rightarrow E$ by $(m,e) \mapsto m \ast e = (me)^{+}$.
If $e \in E$ define $\equiv_{e}$ on $M$ by $m \equiv_{e} n$ iff $em = en$.
You can think of this as being the replacement of an action of $E$ on $M$ which does not work.
It is routine to check (but see the lemma below) that the following axioms hold:
\begin{itemize}
\item[{\rm (MP1).}] The monoid $M$ acts on the set $E$. We denote this action by $(m,e) \mapsto m \ast e$.
\item[{\rm (MP2).}] $m \ast 1 = 1$ for all $m \in M$.
\item[{\rm (MP3).}] $m \ast (ef) = (m \ast e)(m \ast f)$.
\item[{\rm (MP4).}] $\equiv_{e}$ is a right congruence on $M$; the $\equiv_{e}$-class containing the element $m$ will be denoted by $[m]_{e}$.
\item[{\rm (MP5).}] $\equiv_{1}$ is the identity right congruence.
\item[{\rm (MP6).}] If $f \leq e$ then $m \equiv_{e} n$ implies that $m \equiv_{f} n$.
\item[{\rm (MP7).}] If $a \equiv_{e} a'$ then $ma \equiv_{m \ast e} ma'$ for any $m \in M$.
\item[{\rm (MP8).}] If $a \equiv_{e} a'$ then $e(a \ast f) = e(a' \ast f)$ for any $f \in E$.
\end{itemize}
If, now,  $M$ is {\em any} monoid, and $E$ is {\em any} semilattice in which the above axioms hold 
then we will say that we have a {\em matched pair} denoted by $[E|M]$.\footnote{This is in a more general sense than that used by Garner.}

\begin{lemma}
The left restriction monoid $S$ gives rise to the matched pair 
$$[\mathsf{Proj}(S)|\mathsf{Tot}(S)].$$
\end{lemma}
\begin{proof} We have already seen that $M = \mathsf{Tot}(S)$ is a monoid and $E = \mathsf{Proj}(S)$ is 
a meet semilattice with a top element from the definition of left restriction monoid. 
Axiom (MP1) holds: use (LR5) to establish the associativity of the action.
Axiom (MP2) holds from the very definition of $M$.
Axiom (MP3) holds: use (LR2), (LR5) and (LR6).
Axiom (MP4) holds: this is immediate from the definition.
Axiom (MP5) holds: immediate from the properties of the identity.
Axiom (MP6) holds: immediate from $f \leq e$ and properties of the binary operation in $M$.
Axiom (MP7) holds: use (LR6).
Axiom (MP8) holds: use (LR2) and (LR5).
\end{proof}

Let $[E|M]$ and $[E'|M']$ be matched pairs.
A {\em homomorphism} from $[E|M]$ to $[E'|M']$ is a pair of maps $(\alpha, \beta)$
where both $\alpha \colon M \rightarrow M'$ 
and
$\beta \colon E \rightarrow E'$ are homomorphisms,
such that $\beta (m \ast e) = \alpha (m) \ast \beta (e)$, and $m \equiv_{e} n$ implies $\alpha (m) \equiv_{\beta (e)} \alpha (n)$.
In this way, we can see that matched pairs form a category.
If they are {\em isomorphic}, we write  $[E|M] \cong [E'|M']$.

We say that a left restriction monoid $S$ is {\em factorizable} if for each $s \in S$ there is a total element $\hat{s} \in \mathsf{Tot}(S)$ such that $s \leq \hat{s}$.

\begin{remark}{\em Let $S$ be {\em any} left restriction monoid.
Put $S' = \mathsf{Tot}(S)^{\downarrow}$, the elements of $S$ below a total element in the natural partial order.
Then $S'$ is a factorizable left restriction monoid, and it contains all the projections of $S$. 
This example shows that factorizable left restriction monoids are easy to find.}
\end{remark}


%%%%%%%%%%%%%%%%%%%%%%%%%%%%%%%%%%%%%%%%%%%%%%%%%%%%%%%%%%%%%%%%%%%%%%%%%%%%%%%%%%%%%%%%%%%%%%%%%%%%%%%%%%%%%%%
What follows is a mild generalization of \cite[Proposition 3.3, Proposition 3.4, Theorem 3.5]{G2025}.

\begin{theorem}\mbox{}
\begin{enumerate} 

\item Let $[E|M]$ be a matched pair. Then we can construct a factorizable left restriction monoid $S = S[E|M]$
such that $[E|M] \cong [\mathsf{Proj}(S)|\mathsf{Tot}(S)]$.

\item Every factorizable left restriction monoid is isomorphic to a monoid of the form  $S[E|M]$.

\item The category of factorizable left restriction monoids is equivalent to the category of matched pairs. 

\end{enumerate}
\end{theorem}
\begin{proof}
(1) We shall only sketch the proof since it is similar to the one given by Garner in \cite{G2025}.
The elements of $S$ are ordered pairs $(e, [m]_{e})$ where $e \in E$.
Define
$$(e,[a]_{e})(f,[b]_{f}) = (e(a \ast f), [ab]_{e(a \ast f)}).$$
This is well-defined by (MP4), (MP6), (MP7) and (MP8).
We therefore have a binary operation on $S$.
This is associative by (MP1) and (MP3) and has an identity $(1,[1]_{1})$ by (MP1) and (MP2).
It is therefore a monoid.
Define 
$$(e,[a]_{e})^{+} = (e,[1]_{e}).$$
It is now routine to check that $S$ is a left restriction monoid.
The total elements are those of the form $(1,[a]_{1})$.
The projections are the elements of the form $(e,[1]_{e})$.
We have that $(e,[a]_{e}) \leq (f,[b]_{f})$ iff $e \leq f$ and $a \equiv_{e} b$.
It follows that $S$ is factorizable since $(e,[a]_{e}) \leq (1,[a]_{1})$.
Define $\alpha \colon M \rightarrow \mathsf{Tot}(S)$ by $a \mapsto (1,[a]_{1})$,
a bijection by (MP5),
and define $\beta \colon E \rightarrow \mathsf{Proj}(S)$ by $e \mapsto (e,[1]_{e})$, also a bijection.
It is now routine to check that $(\alpha, \beta)$ defines an isomorphism $[E|M] \cong [\mathsf{Proj}(S)|\mathsf{Tot}(S)]$.

(2) Let $S$ be a factorizable left restriction monoid.
Then $[\mathsf{Proj}(S)|\mathsf{Tot}(S)]$ is a matched pair, and so we may construct a factorizable left restriction monoid
$\hat{S} = S[\mathsf{Proj}(S)|\mathsf{Tot}(S)]$ as in (1) above.
It remains to show that $S$ and $\hat{S}$ are isomorphic.
It is now that we use the fact that $S$ is factorizable.
Let $a \in S$.
Then, by assumption, $a = em$ where $e$ is a projection and $m$ is a total element.
If we apply axiom (LR5), then we deduce that $e = a^{+}$.
Although the projection $e$ is uniquely defined,
the element $m$ is not uniquely defined, but if $n$ is any total element above $a$ we must always have $a^{+}m = a^{+}n$
It follows that the map $\theta \colon S \rightarrow \hat{S}$, given by $a \mapsto (a^{+}, [m]_{a^{+}})$, where $m$ is any total element above $a$,
is well-defined.
It is clearly bijective, and an application of axioms (LR5) and (LR6) shows that it is a homomorphism of left restriction monoids.

(3) Let $\theta \colon S \rightarrow T$ be a homomorphism between two factorizable left restriction monoids.
Then it is routine to show that $(\theta | \mathsf{Tot}(S), \theta | \mathsf{Proj}(S))$ is a homomorphism of matched pairs.
Suppose that $(\alpha, \beta)$ is a homomorphism from the matched pair  $[E|M]$ to  the matched pair $[E'|M']$.
Define $\theta \colon S[E|M] \rightarrow S[E'|M']$ by 
$$(e, [a]_{e}) \mapsto (\beta (e), [\alpha (a)]_{\beta (e)}).$$
This is a homomorphism of left restriction monoids. 
The proof of the equivalence of categories now follows from the above and parts (1) and (2).
\end{proof}


The above theorem is a generalization of \cite{EF}, since right congruences on groups are determined by subgroups.




%%%%%%%%%%%%%%%%%%%%%%%%%%%%%%%%%%%%%%%%%%%%%%%%%%%%%%%%%%%%%%%%%%%%%%%%%%%%%%%%%%%%%%%%%%%%%%%%%%%%%%%%%%%%%%%%%%%%%%%%%%%%%%%%%%%%%%%%%%%%%%%%%%%%%%%%%%%%%%%%%%%%%%%%%%%%%%%
\section{Supported actions}

In this section, we shall define a class of actions by left restriction monoids $S$
and describe them in terms of `actions' by the matched pair $[\mathsf{Proj}(S)|\mathsf{Tot}(S)]$.

Let $S$ be a left restriction monoid.
A left action of $S$ on a set $X$, denoted by $(s,x) \mapsto s \cdot x$, is said to be {\em supported}
if there is a function $p \colon X \rightarrow \mathsf{Proj}(S)$, called the {\em support}, such that the following two axioms hold:
\begin{itemize}
\item[{\rm (E1).}] $p(x) \cdot x = x$. 
\item[{\rm (E2).}] $p(s \cdot x) = (sp(x))^{+}$.
\end{itemize}
This definition is modelled on the definition of \'etale inverse semigroup action
to be found in \cite{F2010} and \cite{S2011}. 
We prefer the term `supported' since the word `\'etale' seems to be overused.
In this, we have followed \cite{LR2021}. 
If the left restriction monoid $S$ admits a supported action on the set $X$ with support $p$ then we write $(S,X,p)$.

\begin{examples}\label{exs:examples}
{\em Here are two examples of supported actions.
Let $S$ be a left restriction monoid.
\begin{enumerate}
\item  Put $E = \mathsf{Proj}(S)$.
Define $p \colon E \rightarrow \mathsf{Proj}(S)$ to be the identity function.
Define $S \times E \rightarrow E$ by $(s,e) \mapsto s \cdot e = (se)^{+}$.
It is easy to check that this is a supported action.

\item If $e$ is any projection, then $(S,Se,p_{e})$ where $p_{e}(x) = x^{+}$
is a supported action when the action is defined via multiplication.
\end{enumerate}
}
\end{examples}

Just as in the inverse case, we can define an order $\leq$ on $X$ by $x \leq y$ iff $x = p(x) \cdot y$.
It is easy to check that $x \leq y$ iff $x = e \cdot y$ where $e$ is some projection.

If $x,y \in X$ we write $x \approx y$ if $p(x) \cdot y = p(y) \cdot x$.
Observe that if $x,y \leq z$ then $x \approx y$.

For any projection $e$, we write $X_{e}$ for the set of all elements $x \in X$ such that $p(x) = e$.
Particularly significant to us will be the set $X_{1}$.

\begin{remark}{\em Suppose that we have a supported action $(S,X,p)$.
Then $X$ carries the structure of a presheaf of sets over the semilattice $\mathsf{Proj}(S)$ as follows:
if $f$ and $e$ are projections such that $f \leq e$ define $\phi^{e}_{f} \colon X_{e} \rightarrow X_{f}$ by $x \mapsto f \cdot x$;
observe that $\phi^{e}_{e}$ is the identity function on $X_{e}$ and if $g \leq f \leq e$ then $\phi^{e}_{g} = \phi^{f}_{g}\phi^{e}_{f}$.}
\end{remark}

We say that $X$ is {\em factorizable} if for each $y \in Y$ there exists $x \in X_{1}$
such that $y \leq x$.

Let $(S,X,p)$ and $(S,X',p')$ be supported actions of a left restriction monoid $S$.
A {\em homomorphism} of such actions is a function $\theta \colon X \rightarrow X'$ 
which satisfies two conditions:
\begin{enumerate}
\item $\theta$ is $S$-equivariant. 
\item $p' \theta = p$ 
\end{enumerate}
Denote by $\mbox{hom}(X,X')$ the set of all such homomorphisms.

Fix a left restriction monoid $S$.
We then obtain a category of all supported actions of $S$ which we denote by $S$-{\bf Sup}.
We shall first of all say something about this category.
In what follows, 
we use the definition of Cartesian closed category to be found in, for example, \cite{BW}: 
thus we require a terminal object,
binary products and exponentials.

\begin{theorem}\label{them:CCV1} Let $S$ be a fixed left restriction monoid.
Put $\mathscr{C} = S$-{\bf Sup}.
\begin{enumerate}
\item $\mathscr{C}$ is Cartesian closed.
\item The subcategory of $\mathscr{C}$ consisting of those actions with surjective support is Cartesian closed.
\item The subcategory of $\mathscr{C}$ consisting of factorizable actions is Cartesian closed.
\end{enumerate}
\end{theorem}
\begin{proof} We prove (1) which is the main result, and indicate how (2) and (3) are proved as we go along.

We show first that the category $\mathscr{C}$ has a terminal object.
Put $E = \mathsf{Proj}(S)$.
Define $p \colon E \rightarrow \mathsf{Proj}(S)$ to be the identity function.
Define $S \times E \rightarrow E$ by $(s,e) \mapsto s \cdot e = (se)^{+}$.
By part (1) of Examples~\ref{exs:examples} this is a supported action.
Let $(S,X,q)$ be any supported action.
Define $\theta \colon X \rightarrow E$ by $\theta (s) = q(x)$.
It is easy to check that $\theta$ is a homomorphism of supported $S$-actions.
It is also the unique such homomorphism.
We have therefore proved that $(S,E,p)$ is the terminal object in $\mathscr{C}$. 
Observe that $p$ is always surjective and the action is always factorizable.

We now show that the category $\mathscr{C}$ has binary products.
Let $(S,X,p)$ and $(S,Y,q)$ be supported actions of $S$.
Define $Z$ to the set of all ordered pairs $(x,y)$ where $p(x) = q(y)$. 
Define $r \colon Z \rightarrow E(S)$ by $r(x,y) = p(x) = q(y)$.
There is an action of $S$ on $Z$ given by $s \cdot (x,y) = (s \cdot x, s \cdot y)$.
This is well-defined and gives rise to a supported action $(S,Z,r)$.
In fact, this action is the product of $(S,X,p)$ and $(S,Y,q)$ in the category $\mathscr{C}$.
For clarity, we shall denote this product action by $X \,\Box \, Y$.
Observe that if $p$ and $q$ are surjective then $r$ is surjective.
Suppose now that both $X$ and $Y$ are factorizable.
Let $(x,y) \in X \, \Box \, Y$.
By assumption, $x = p(x) \cdot x'$ and $y = q(y) \cdot y'$  where $x' \in X_{1}$ and $y' \in Y_{1}$.
Observe that $r(x',y') = 1$ and $(x,y) = r(x,y) \cdot (x',y')$.
Thus $X \, \Box \, Y$ is factorizable.

Given supported actions $(S,X,p)$ and $(S,Y,q)$ and using part (2) of Examples~\ref{exs:examples},
define
$$Y^{X} = \bigsqcup_{e \in \mathsf{Proj}(S)} \mbox{hom}(Se \, \Box \, X, Y),$$ 
and define $\pi \colon Y^{X} \rightarrow \mathsf{Proj}(S)$ by $\pi (\theta) = e$ if $\theta \in \mbox{hom}(Se \, \Box \, X,Y)$.
Observe that $\pi$ is always surjective.
We now define a function $S \times Y^{X} \rightarrow Y^{X}$ as follows: 
if $\theta \colon Se\, \Box \,X \rightarrow y$ then define
$$(a \cdot \theta)(t,x') = \theta (ta, (ta)^{+} \cdot x').$$
This is well-defined and $a \cdot \theta \in \mbox{hom}(Se\, \Box X,Y)$.
It can be checked that $(S,Y^{X},\pi)$ is a supported action.
Let $(S,Z,r)$  be any supported action.
Suppose that $g \colon Z\, \Box \, X \rightarrow Y$ is a homomorphism.
Define $\hat{g} \colon Z \rightarrow Y^{X}$ as follows:
we have that $\hat{g}(z) \colon Sr(z) \, \Box \, X \rightarrow Y$ defined by
$$\hat{g}(z)(a,x) = g(a \cdot z, (ar(z))^{+} \cdot x).$$
It can be checked that
$$\hat{g}(m \cdot z) = (m \cdot \hat{g})(z).$$
Define $\mathbf{eval} \colon Y^{X} \, \Box \, X \rightarrow Y$ as follows:
if $\pi (\theta) = e$ 
then $\mathbf{eval}(\theta, x) = \theta (e,x)$.
It is easy to check that $\mathbf{eval}(\hat{g}(z),x) = g(z,x)$.

It only remains to prove the `uniqueness' part of the definition of Cartesian closed.
Suppose that $\bar{g} \colon Z \rightarrow Y^{X}$ is a homomorphism
such that $\bar{g}(z)(r(z),x) = \hat{g}(z)(r(z), x)$.
We show that $\bar{g} = \hat{g}$.
To do this, we show that $\bar{g}(a,x) = \hat{g}(z)(a,x)$
where $a \in Sr(z)$ and $a^{+} = p(x)$.
Observe that 
$$\hat{g}(z)(a,x) = (a \cdot \hat{g})(z)(a^{+},  x) = \hat{g}(a \cdot z)(a^{+}, x) = \hat{g}(a \cdot z)(r(a \cdot z), (ar(z)))^{+} \cdot x).$$
But, by assumption, 
$$\hat{g}(a \cdot z)(r(a \cdot z), (ar(z))^{+} \cdot x) = \bar{g}(a \cdot z)(r(a \cdot z), (ar(z)))^{+} \cdot x).$$
We also have that 
$$\bar{g}(a \cdot z)(r(a \cdot z), (ar(z)))^{+} \cdot x) = (a \cdot \bar{g})(z)(r(a \cdot z), (ar(z)))^{+} \cdot x).$$ 
But
$$(a \cdot \bar{g})(z)(r(a \cdot z), (ar(z)))^{+} \cdot x)
= 
\bar{g}(z)(r(a \cdot z)a, (ar(z))^{+} \cdot x) = \bar{g}(z)(a,x).$$
We prove that $Y^{X}$ is always factorizable.
Let $\theta \colon Se\, \Box \, X \rightarrow Y$ be any homomomorphism.
Define $\phi$ as follows:
if $(s,x) \in S \, \Box \, X$ then $\phi (s,x) = \theta (se, (se)^{+} \cdot x)$.
This is a well-defined function and it is a homomorphism by axiom (LR6).
Observe that $e \cdot \phi$ is a homomorphism from $Se \, \Box \, X$ to $Y$.
It is easy to check that $\theta = e \cdot \phi$.
\end{proof}

\begin{example}{\em  Let $M$ be any monoid with identity 1.
Then $M$ is a left restriction monoid if we define $m^{+} = 1$ for all $m \in M$.
In this case, the set of projections is just $\{1\}$.
Supported actions are simply monoid actions.
The product action of two actions has as underlying set the product of the two sets.
The exponential reduces to $Y^{X} = \mbox{hom}(M \times X,Y)$ in this case.}
\end{example}


In the light of the above theorem, the following is a natural question.

\begin{question}
{\em Which Cartesian closed categories are equivalent to those of the form $S$-{\bf Sup} for some left restriction monoid $S$?}
\end{question}

%%%%%%%%%%%%%%%%%%%%%%%%%%%%%%%%%%%%%%%%%%%%%%%%%%%%%%%%%%%%%%%%%%%%%%%%%%%%%%%%%%%%%%%%%%%%%%%%%%%%%%%%%%%%%%%%%%%%%%%%%%%%%%%%%%%%%%%%%%%%
Let $S$ be a left restriction monoid.
We denote the category of factorizable supported actions of $S$ by $S$-{\bf FSup}.
We shall obtain another description of this category.
We shall carry forward any notation introduced in Section~2.

Let $(S,X,p)$ be a factorizable (although, initially, this will play no role) supported action.
Put $M = \mathsf{Tot}(S)$, a monoid, and $E = \mathsf{Proj}(S)$, a meet semilattice with top element, and $Y = X_{1}$.
We therefore have a matched pair $[E|M]$.
The action of $S$ on $X$ leads to an action of $M$ on $Y$, as you can check.
{\em However, $E$ does not act of $Y$}, and this is the key point.
Following Garner, we define the relation $\equiv_{e}$ on $Y$ for each $e \in E$ by $x \equiv_{e} y$, for $x,y \in Y$, iff $e \cdot x = e \cdot y$.
The equivalence relation $\equiv_{e}$ contains some information about the action of $e$ on $Y$.
The following axioms hold:
\begin{itemize}
\item[{\rm (MPA1).}] The monoid $M$ acts on the set $Y$.
\item[{\rm (MPA2).}] For each $e \in E$ there is an equivalence relation $\equiv_{e}$
\item[{\rm (MPA3).}] $\equiv_{1}$ is the identity equivalence.
\item[{\rm (MPA4).}] If $f \leq e$ then $x \equiv_{e} y$ implies that $x \equiv_{f} y$.
\item[{\rm (MPA5).}] If $m \equiv_{e} n$ in $M$ and $x \in Y$ then $m \cdot x \equiv_{e} n \cdot x$.
\item[{\rm (MPA6).}] If $x \equiv_{e} y$ in $Y$ then for any $m \in M$ we have that $m \cdot x \equiv_{m \ast e} m \cdot y$.
\end{itemize}
{\em Observe that whereas $M$ really does act on the set $Y$, the monoid $E$ does not act on $Y$.
Nevertheless, we shall regard this as an `action' of the matched pair $[E|M]$ on the set $Y$.}
We call the set $Y$ an {\em $[E|M]$-set}.

If $Y$ and $Y'$ are $[E|M]$-sets, then a {\em homomorphism} from $Y$ to $Y'$ is a function $\alpha \colon Y \rightarrow Y'$
such that $\alpha$ is a map of $M$-actions and if $x \equiv_{e} y$ then $\alpha(x) \equiv_{e} \alpha (y)$.
We may therefore form the category of $[E|M]$-sets which we denote by $[E|M]$-{\bf Set}.

\begin{theorem}\label{them:mary} Let $S$ be a left restriction monoid.  
The category $S$-{\bf FSup} is equivalent to the category $[\mathsf{Proj}(S)|\mathsf{Tot}(S)]$-{\bf Set}.
\end{theorem}
\begin{proof} 
We first define a functor from $S$-{\bf FSup} to $[\mathsf{Proj}(S)|\mathsf{Tot}(S)]$-{\bf Set}.
We start with a supported action $(S,X,p)$.
As above, we have a $[\mathsf{Proj}(S)|\mathsf{Tot}(S)]$-set $X_{1}$.
Suppose we also have a supported action $(S,X',p')$ and a homomorphism of supported actions $\theta \colon X \rightarrow X'$.
Put $\theta' = \theta | X_{1}$.
Suppose that $p(y) = 1$ then $p' (\theta (y)) = 1$.
It follows that $\theta' \colon X_{1} \rightarrow X'_{1}$ is a well-defined function.
Suppose that $x \equiv_{e} x'$ in $X_{1}$.
Then, by definition, $e \cdot x = e \cdot x'$.
It follows that $e \cdot \theta (x) = e \cdot \theta (x')$.
Thus $\theta'$ is a homomorphism of $[\mathsf{Proj}(S)|\mathsf{Tot}(S)]$-sets.
We have therefore defined our functor.

We now define a functor from $[\mathsf{Proj}(S)|\mathsf{Tot}(S)]$-{\bf Set} to 
$S$-{\bf FSup}.
We start with a $[\mathsf{Proj}(S)|\mathsf{Tot}(S)]$-set $Y$.
Define
$$X = \bigcup_{e \in \mathsf{Proj}(S)} Y/\equiv_{e}$$
and
$p \colon X \rightarrow \mathsf{Proj}(S)$
by $p([x]_{e}) = e$.
We now define what will turn out to be a supported action of $S$ on $X$.
Let $s = s^{+}\hat{s}$ where $\hat{s} \in \mathsf{Tot}(S)$.
Define 
$$s \bullet [x]_{e} = [\hat{s} \cdot x]_{(se)^{+}}.$$
We show that this is well-defined.
Let $[x]_{e} = [x']_{e}$ and 
$s = s^{+}\hat{s} = s^{+}\hat{t}$, where $\hat{s}, \hat{t} \in \mathsf{Tot}(S)$.
We prove that
$[\hat{s} \cdot x]_{(se)^{+}} = [\hat{t} \cdot x']_{(se)^{+}}$.
From $\hat{s} \equiv_{s^{+}} \hat{t}$ we get that $\hat{s} \cdot x \equiv_{s^{+}} \hat{t} \cdot x$
by axiom (MPA5).
From $x \equiv_{e} x'$ we get $\hat{t} \cdot x \equiv_{\hat{t} \ast e} \hat{t} \cdot x'$
by axiom (MPA6).
We now apply axiom (MPA4) twice.
From $(se)^{+} \leq s^{+}$,
we get $\hat{s} \cdot x \equiv_{(se)^{+}} \hat{t} \cdot x$,
and from $(se)^{+} \leq (\hat{t}e)^{+}$,
we get $\hat{t} \cdot x \equiv_{(se)^{+}} \hat{t} \cdot x'$.
It follows that	
$$\hat{s} \cdot x \equiv_{(\hat{s}e)^{+}} \hat{t} \cdot x'.$$
It is now easy to show that $(S,X,p)$ is a supported action.
Observe that this is actually factorizable: 
let $[x]_{e}$ be any element of $X$.
Then $e \bullet [x]_{1} = [x]_{e}$.
We now turn to homomorphisms.
Suppose that $\alpha \colon Y \rightarrow Y'$ is a homomorphism of 
$[\mathsf{Proj}(S)|\mathsf{Tot}(S)]$-sets.
Let $(S,X,p)$ and $(S,X',p')$ be the corresponding supported actions.
Then we may define a homomorphism $\theta$ of supported actions from $X$ to $X'$
by
$$\theta ([y]_{e}) = [\alpha (y)]_{e}.$$
We have therefore defined a functor from the category 
$[\mathsf{Proj}(S)|\mathsf{Tot}(S)]$-{\bf Set} to 
the category $S$-{\bf FSup}

To prove that these two functors define an equivalence of categories,
we have to prove two things; we do this below.

Suppose first that $(S,X,p)$ is a supported action.
Put $X' = \bigcup_{e \in \mathsf{Proj}(S)} X_{1}/\equiv_{e}$.
Define $\theta \colon X \rightarrow X'$ by $\theta (x) = [y]_{p(x)}$
where $p(y) = 1$ and $x \leq y$.
This is an isomorphism of supported actions.

Let $Y$ be a $[\mathsf{Proj}(S)|\mathsf{Tot}(S)]$-set.
The associated supported action set is $X = \bigcup_{e \in \mathsf{Proj}(S)} x/\equiv_{e}$
and the associated support is $p([x]_{e}) = e$.
Thus the set $X_{1}$ consists of singleton sets $\{y\}$ (by axiom (MPA3)) where $y \in Y$.
There is an obvious isomorphism.  
\end{proof}

The following will be needed later.
The proof is routine.

\begin{lemma}\label{lem:agatha}
Let $S$ be a left restriction monoid.
Let $Y$ be a $[\mathsf{Proj}(S)|\mathsf{Tot}(S)]$-set.
The semigroup $S$ acts on the set $X$ where
$$X = \bigcup_{e \in \mathsf{Proj}(S)} Y/\equiv_{e}$$
where
$$[x']_{f} \leq [x]_{e} \text{ iff } f \leq e \text{ and } x' \equiv_{f} x.$$ 
\end{lemma}



%%%%%%%%%%%%%%%%%%%%%%%%%%%%%%%%%%%%%%%%%%%%%%%%%%%%%%%%%%%%%%%%%%%%%%%%%%%%%%%%%%%%%%%%%%%%%%%%%%%%%%%%%%%%%%%%%%%%%%%
\section{The Boolean case}

In this section, we shall refine what we have shown in the previous two sections
and make contact with Garner's work.
To this end, we shall not work with arbitrary left restriction monoids but those which are `Boolean'.
Our first goal, then, is to define what this means.


Let $S$ be a left restriction monoid.
Define $a \sim_{r} b$, on $S$, and say that $a$ is {\em right-compatible} with $b$,tube
if $a^{+}b = b^{+}a$.
Being right-compatible is a necessary condition for a pair of elements to have a join.
If $S$ has a zero which is also a projection,
then we say that the elements $s$ and $t$ are {\em right-orthogonal}, denoted by $s \perp t$, if $s^{+}$ and $t^{+}$ are orthogonal (meaning $s^{+}t^{+} = 0$).
Right-orthogonal elements are right-compatible.
We say that a left restriction monoid $S$ is {\em Boolean} if it satisfies the following three axioms:
\begin{itemize}
\item[{\rm (B1).}] $\mathsf{Proj}(S)$ is a Boolean algebra with respect to the natural partial order
with the zero of the semigroup being the zero of the Boolean algebra, and the identity of the monoid being the
one of the Boolean algebra. 
\item[{\rm (B2).}] Right-compatible pairs of elements have joins. We say that {\em binary joins exist}.
\item[{\rm (B3).}] Multiplication distributes over binary joins from the left and from the right. 
\end{itemize}
Despite appearances, our definition of Boolean left restriction monoid agrees with the definition
of Boolean restriction monoid given as \cite[Definition 3.2]{G2025}.
Homomorphisms of Boolean left restriction monoids are required to preserve binary joins.
The left-right dual of the following result is proved in \cite[Lemma 3.1]{L2025}.
It will be used many times.

\begin{lemma}\label{lem:induction} In any Boolean left restriction monoid,
if $a = b \vee c$ exists then $a^{+} = b^{+} \vee c^{+}$.
\end{lemma}

The following will be needed later, but answers an obvious question.

\begin{lemma}\label{lem:people} Let $S$ be a Boolean left restriction monoid.
Let $s \in S$ and $e$ be any projection.
Then $s\,\bar{e} = \overline{(se)^{+}}\,s$.
\end{lemma}
\begin{proof} We have that $1 = e \vee \bar{e}$ and so $s = se \vee s \bar{e}$.
Similarly, $1 = (se)^{+} \vee \overline{(se)^{+}}$
and so $s = (se)^{+}s \vee \overline{(se)^{+}}s = se \vee \overline{(se)^{+}}s$.
We have therefore shown that $s\bar{e} \sim_{r} \overline{(se)^{+}}s$, since both elements are beneath $s$.
The result will follow if we can show that $(s\bar{e})^{+} = (\overline{(se)^{+}}s)^{+}$.
However, this follows from standard Boolean algebra
given that $s^{+} = (se)^{+} \vee (s\bar{e})^{+}$, by Lemma~\ref{lem:induction}, is an orthogonal join
and $s^{+} = (se)^{+} \vee (\overline{(se)^{+}}s)^{+}$, by Lemma~\ref{lem:induction}, is also an orthogonal join.
\end{proof}

The next result connects what we are doing here with Section~2.

\begin{lemma} 
Every Boolean left restriction monoid is factorizable.
\end{lemma}
\begin{proof} Let $S$ be a Boolean left restriction monoid.
Let $s \in S$.
Then the elements $s$ and $\overline{s^{+}}$ are right-orthogonal.
If we put $\hat{s} = s \vee \overline{s^{+}}$ then we obtain a total element, by Lemma~\ref{lem:induction}, that is above $s$ in the natural partial order.
We may therefore write $s = s^{+}\hat{s}$.
It follows that Boolean left restriction monoids are automatically factorizable.
\end{proof}

%%%%%%%%%%%%%%%%%%%%%%%%%%%%%%%%%%%%%%%%%%%%%%%%%%%%%%%%%%%%%%%%%%%%%%%%%%%%%%%%%%%%%%%%%%%%%%%%%%%%%%%%%%%%%%%%%%%%%%%%%%%%%%%%%%%%%%%%%%
\begin{center}
{\bf The construction of Boolean left restriction monoids from matched pairs}
\end{center}

We shall now strengthen our definition of matched pairs to the case where $E = B$ is actually a Boolean algebra.
This will then agree with the definition of matched pair that Garner gave.
Let $M$ be a monoid and $B$ a Boolean algebra.
We say that $[B|M]$ is a {\em matched pair} if the following axioms hold:
\begin{itemize}
\item[{\rm (MP1).}] The monoid $M$ acts on the set $B$.
\item[{\rm (MP2).}] $m \ast 1 = 1$ for all $m \in M$.
\item[{\rm (MP3).}] $m \ast (ef) = (m \ast e)(m \ast f)$.
\item[{\rm (MP4).}] $\equiv_{e}$ is a right congruence on $M$. The $\equiv_{e}$-class containing the element $m$ will be denoted by $[m]_{e}$.
\item[{\rm (MP5).}] $\equiv_{1}$ is the identity right congruence.
\item[{\rm (MP6).}] If $f \leq e$ then $m \equiv_{e} n$ implies that $m \equiv_{f} n$.
\item[{\rm (MP7).}] If $a \equiv_{e} a'$ then $ma \equiv_{m \ast e} ma'$ for any $m \in M$.
\item[{\rm (MP8).}] If $a \equiv_{e} a'$ then $e(a \ast f) = e(a' \ast f)$ for any $f \in E$.
\item[{\rm (MP9).}] $\equiv_{0}$ is the universal right congruence.
\item[{\rm (MP10).}] $m \ast (e + f) = m \ast e + m \ast f$.
\item[{\rm (MP11).}] if $m \equiv_{e} n$ and $m \equiv_{f} n$ then $m \equiv_{e+f} n$. 
\item[{\rm (MP12).}] Given any\ $m,n \in M$ and any $e \in B$, then there exists a $p \in M$ such that $p \equiv_{e} m$ and $p \equiv_{\bar{e}} n$. 
\end{itemize}
These are the same axioms as in Section~2 but augmented by the axioms (MP9)--(MP12) to take account of the extra structure in a Boolean algebra
compared with a semilattice.


Our first result connects Boolean left restriction monoids with matched pairs.


\begin{lemma}\label{lem:anne} Let $S$ be a Boolean left restriction monoid.
Then $[\mathsf{Proj}(S)|\mathsf{Tot}(S)]$ is a matched pair (in the above sense).
\end{lemma}
\begin{proof} We know from Section~2, that the axioms (MP1)--(MP8) all hold.
It is immediate that axioms (MP9) and (MP10) hold.
To see that (MP11) holds, suppose that $m \equiv_{e} n$ and $m \equiv_{f} n$ both hold.
Then $em = en$ and $fm = fn$.
It is easy to check that $(e + f)m = (e + f)n$, and the result follows.
For axiom (MP12), put $p = em \vee \bar{e}n$.
This makes sense because $em$ and $\bar{e}n$ are right-orthogonal.
\end{proof}

The above result raises the question of whether every Boolean left restriction monoid
arises from such a matched pair. 
The answer is in the affirmative and was answered first in \cite{G2025}.
We shall need the follwing lemma.

\begin{lemma}\label{lem:amelia} Let $[B|M]$ be a matched pair where $B$ is a Boolean algebra.
Suppose that $m \equiv_{ef} n$. Then there exists a $p \in M$ such that
$p \equiv_{e} m$ and $p \equiv_{f} n$.
Furthermore, if $p' \in M$ is any element such that 
$p' \equiv_{e} m$ and $p' \equiv_{f} n$ then
$p \equiv_{e + f} p'$.
\end{lemma}
\begin{proof} By axiom (MP12), there exists $p \in M$ such that $p \equiv_{e} m$ and $p \equiv_{\bar{e}} n$.
Observe that $f = ef + \bar{e}f$.
By axiom (MP6), we have that $p \equiv_{ef} m$ and $p \equiv{\bar{e}f}$.
But $m \equiv_{ef} n$ so that $p \equiv{ef} n$.
Thus $p \equiv_{ef} m$ and $p \equiv_{\bar{e}f} n$ and so by axiom (MP11),
we have that $p \equiv_{f} n$.
Suppose that $p' \equiv_{e} m$ and $p' \equiv_{f} n$.
Then by axiom (MP6), we have that $p' \equiv_{e} m$ and $p' \equiv_{f} n$.
Thus $p' \equiv_{e} p$ and $p' \equiv_{f} p$.
It follows from axiom (MP11),  that $p' \equiv_{e + f} p$ 
\end{proof}


We now have the following result.
This is proved in \cite{G2025} so we shall just sketch it here (though we take a slightly different tack).


\begin{theorem} Let $M$ be a monoid and $B$ a Boolean algebra. 
\begin{enumerate} 

\item Let $[B|M]$ be a matched pair. Then we can construct a Boolean left restriction monoid $S = S[B|M]$.
In addition, $[B|M] \cong [\mathsf{Proj}(S)|\mathsf{Tot}(S)]$.

\item Every Boolean left restriction monoid is isomorphic to a monoid of the form  $S[B|M]$.

\item The category of Boolean left restriction monoids is equivalent to the category of matched pairs. 

\end{enumerate}
\end{theorem}
\begin{proof} Suppose that $(e,[a]_{e}) \sim_{r} (f,[b]_{f})$.
Then $a \equiv_{ef} b$.
By Lemma~\ref{lem:amelia}, there is an element $p \in M$ such that $p \equiv_{e} m$ and $p \equiv_{f} n$.
Consider the element $(e+f, [p]_{e+f})$.
It is now easy to show that
$$(e,[a]_{e}) \vee (f,[b]_{f}) = (e+f, [p]_{e+f}).$$
The remainder of the proof follows
\cite[Proposition 3.4, Theorem 3.5]{G2025}.
\end{proof}


%%%%%%%%%%%%%%%%%%%%%%%%%%%%%%%%%%%%%%%%%%%%%%%%%%%%%%%%%%%%%%%%%%%%%%%%%%%%%%%%%%%%%%%%%%%%%%%%%%%%%%%%%%%%%%%%%%%%%%%%%%%%%%%%%%%%%%%%%%%%%%
\begin{center}
{\bf Supported actions of  Boolean left restriction monoids and `actions' of matched pairs}
\end{center}


We need to modify what we mean by a supported action in the case where we are acting by Boolean left restriction monoids.

If $S$ is a Boolean left restriction monoid, then a supported action of $S$ on a set $X \neq \varnothing$ is required to satisfy the following axioms {\em in addition to (E1) and (E2)}:
\begin{itemize}
\item[{\rm (E3).}] The poset $(X,\leq)$ has a mimimum element $z$ such that $s \cdot z = z$ for all $s \in S$.
\item[{\rm (E4).}] $0 \cdot x = z$ for all $x \in X$. 
\item[{\rm (E5).}] If $x \approx y$ then $x \vee y$ exists in $X$ and $p(x \vee y) = p(x) \vee p(y)$.
\item[{\rm (E6).}] If $x \approx y$ then $s \cdot (x \vee y) = s \cdot x \vee s \cdot y$.
\item[{\rm (E7).}] If $s \sim_{r} t$ then $(s \vee t) \cdot x = s \cdot x \vee t \cdot x$.
\end{itemize}

\begin{remark}{\em Observe that axiom (E6) makes sense since if $x \approx y$ then for any $s$ we have that
$s \cdot x \approx s \cdot y$. In addition, axiom (E7) makes sense since if $s \sim_{r} t$ then $s \cdot x \approx t \cdot x$
for any elements $s,t \in S$.}
\end{remark}

\begin{remark}{\em 
Supported actions of Boolean inverse monoids are defined in the same way except that in axiom (E7)
the condition $s \sim_{r} t$ is replaced by $s \sim t$.}
\end{remark}

\begin{remark}{\em Supported actions of Boolean right restriction monoids
will always be assumed to satisfy the additional axioms (E3)--(E7).}
\end{remark}

Let $(S,X)$ and $(S,X')$ be supported actions of the Boolean left restriction monoid $S$.
A {\em homomorphism} of such actions is a function $\theta \colon X \rightarrow X'$ 
which satisfies the two conditions for a homomorphism of supported actions and the two conditions below:
\begin{enumerate}
\item The smallest element of $X$ is mapped to the smallest element of $X'$.
\item Binary joins are preserved. 
\end{enumerate}

This means that for a fixed Boolean left restriction monoid $S$,
we can define the category whose objects are the supported actions $(S,X,p)$.
This category will also be denoted by $S$-{\bf Supp} though we stress 
that we assume that acxioms (E3)--(E7) also hold when $X$ is non-empty.


\begin{theorem}\label{them:CCV2} Let $S$ be a fixed Boolean left restriction monoid.
Then $\mathscr{C} = S$-{\bf Sup} is Cartesian closed.
\end{theorem}
\begin{proof} This is a refinement of Theorem~\ref{them:CCV1}.  
The terminal object belongs to this category.
If $(S,X,p)$ and $(S,X',p')$ are supported actions
then so too is $(S,X \,\Box \,X',r)$.
It remains to look at exponentials.
Observe that if $e \leq f$ then $Se \, \Box \, X \subseteq Sf \, \Box \, x$.
We may therefore regard $Y^{X}$ as a set of partial functions of $S \, \Box \, X$.
Observe that $\theta \leq \phi$ iff $\theta \subseteq \phi$
and that $\theta \approx \phi$ when $\theta$ and $\phi$ agree on their intersection.
Therefore $\theta \vee \phi = \theta \cup \phi$, when defined.
the zero element is $\zeta \colon \{ 0\} \, \Box \, X \rightarrow Y$ given by
$\zeta \colon \{(0,z)\} \rightarrow \{ z\}$.
\end{proof}


%%%%%%%%%%%%%%%%%%%%%%%%%%%%%%%%%%%%%%%%%%%%%%%%%%%%%%%%%%%%%%%%%%%%%%%%%%%%%%%%%%%%%%%%%%%%%%%%%%%%%%%%%
Let $B$ be a Boolean algebra and $M$ a monoid such that $[B|M$] is a matched pair.
Let $X$ be any set.
We say that $X$ is a {\em $[B|M]$-set} if the following axioms hold:
\begin{itemize}
\item[{\rm (MPA1).}] The monoid $M$ acts on the set $X$.
\item[{\rm (MPA2).}] For each $e \in E$ there is an equivalence relation $\equiv_{e}$ on $X$.
\item[{\rm (MPA3).}] $\equiv_{1}$ is the identity equivalence.
\item[{\rm (MPA4).}] If $f \leq e$ then $x \equiv_{e} y$ implies that $x \equiv_{f} y$.
\item[{\rm (MPA5).}] If $m \equiv_{e} n$ in $M$ and $x \in X$ then $m \cdot x \equiv_{e} n \cdot x$.
\item[{\rm (MPA6).}] If $x \equiv_{e} y$ in $X$ then for any $m \in M$ we have that $m \cdot x \equiv_{m \ast e} m \cdot y$.
\item[{\rm (MPA7).}] $\equiv_{0}$ is the universal equivalence.
\item[{\rm (MPA8).}] if $x \equiv_{e} x'$ and $x \equiv_{f} x'$ then $x \equiv_{e+f} x'$. 
\item[{\rm (MPA9).}] Given any $x,y \in X$ and any $e \in B$, then there exists a $z \in X$ such that $z \equiv_{e} x$ and $z \equiv_{\bar{e}} y$. 
\end{itemize}
These are the same axioms as in Section~3, but augmented by the axioms (MPA7), (MPA8) and (MPA9) to take account of the extra structure in a Boolean algebra.
{\em Homomorphisms} of such $[B|M]$-sets are defined as before.


\begin{lemma} Let $(S,X,p)$ be a factorizable supported action of a Boolean left restriction monoid.
\begin{enumerate}
\item If $X_{1}$ is empty then $X$ is empty.
\item $X_{1}$ is non-empty  iff $p$ is surjective.
\end{enumerate}
\end{lemma}
\begin{proof} (1) Immediate
(2) Suppose that $X_{1}$ is non-empty.
Let $x \in X_{1}$ and let $e$ be any projection.
Then $p(e \cdot x) = e$.
Thus $p$ is surjective.
Conversely, if $p$ is surjective then $X_{1} = p^{-1}(1)$ is automatcally non-empty.
\end{proof}


We now have the following version of Theorem~\ref{them:mary} for Boolean left restriction monoids.

\begin{theorem}\label{them:marytwo} Let $S$ be a Boolean left restriction monoid.  
The category $S$-{\bf FSup} is equivalent to the category $[\mathsf{Proj}(S)|\mathsf{Tot}(S)]$-{\bf Set}.
\end{theorem}
\begin{proof} Put $M = \mathsf{Tot}(S)$, $B = \mathsf{Proj}(S)$. 
We first define a functor from $S$-{\bf FSup} to $[B|M]$-{\bf Set}.
In fact, $[B|M]$ is a matched pair by Lemma~\ref{lem:anne}.
Suppose $(S,X,p)$ is a supported action.
Put $X_{1} = p^{-1}(1)$.
Then $X_{1}$ is a $[B|M]$-set by Theorem~\ref{them:mary}.
It just remains to check that the additional axioms hold.
It is straightforward to show that (MPA7) and (MPA8) hold.
To show that axiom (MPA9) holds,
let $e \in B$ and $x,y \in X_{1}$.
Recall that $x \equiv_{e} y$ iff $e \cdot x = e \cdot y$.
Observe that
$e \cdot x \approx \bar{e} \cdot y$.
Thus $z = e \cdot x \vee \bar{e} \cdot x$ is well-defined.
But $z \equiv_{e} x$ and $z \equiv_{\bar{e}} y$.
Thus (MPA9) really holds.
Let $(S,X,p)$ and $(S,X',p')$ be supported actions
and
let $\theta \colon X \rightarrow X'$ be a homomorphism of supported actions.
The proof that $\theta | X_{1} \colon X_{1} \rightarrow X'_{1}$ defines a homomorphism of $[B|M]$-sets
is now routine.
It is now easy to see that we have defined a functor from the category of supported actions of $S$
to the category {\bf $[B|M]$-sets}.

We now define a functor from $[B|M]$-{\bf Set} to  $S$-{\bf FSup}.
This will follow Theorem~\ref{them:mary}.
Let $Y$ be a $[B|M]$-set.
We assume that $Y$ is non-empty.
As before, 
define
$$X = \bigcup_{e \in \mathsf{Proj}(S)} Y/\equiv_{e}$$
and
$p \colon X \rightarrow \mathsf{Proj}(S)$
by $p([x]_{e}) = e$.
Define 
$$s \bullet [x]_{e} = [\hat{s} \cdot x]_{(se)^{+}}.$$
This defines a supported action of $S$ on $X$, but we now need to show that the axioms
(E3)--(E7) hold for $X$.

(E3) holds: 
Put $z = [y]_{0}$ where $y \in Y$ is any element.
By (MPA7), it does not matter which element $y$ we choose.
From the definition of the action, this really is the minimum element of the poset $X$.
Again from the definition of the action, we have proved that this axiom holds.

(E4) holds: clear.

(E5) holds.
We need a preliminary result first:
if $x \equiv_{ef} y$ then there exists a $w \in X$ such that
$w \equiv_{e} x$ and $w \equiv_{f} y$.
Such a $w$ is unique upto $\equiv_{ef}$-equivalence.
By (MPA9)
there is a $w \in X$ such that $w \equiv_{e} x$ and $w \equiv_{\bar{e}} y$.
In a Boolean algebra, we have that $f = ef \vee \bar{e}f$.
By (MPA4),
we have that $w \equiv_{ef} x$ and $w \equiv_{\bar{e}f} y$. 
It follows that $w \equiv_{ef} y$ and $w \equiv_{\bar{e}f} y$.
By (MPA8), we have that that $w \equiv_{f} y$.
Now, suppose that there exists a $w' \in X$ such that
$w' \equiv_{e} x$ and $w' \equiv_{f} y$.
Then $w \equiv_{e} w'$.
But $ef \leq e$.
Thus by (MPA4),
we have that  $w \equiv_{ef} w'$, as required.

We can now prove that axiom (E5) holds.
Suppose that  $[x]_{e} \approx [y]_{f}$.
We prove that
$[w]_{e + f} = [x]_{e} \vee [y]_{f}$
where $w \in X$ is any element such that 
$w \equiv_{e} x$ and $w \equiv_{f} y$.
Observe that $x \equiv_{ef} y$ means precisely $[x]_{e} \approx [y]_{f}$.
Thus if $[x]_{e} \approx [y]_{f}$, then, using the element $w$ guaranteed above,
we have that $[x]_{e}, [y]_{f} \leq [w]_{e + f}$.
Suppose that $[x]_{e}, [y]_{f} \leq [w']_{i}$.
Then $e \leq i$ and $x \equiv_{e} w'$, and $f \leq i$ and $y \equiv_{f} w'$.
Thus $e + f \leq i$, and $w \equiv_{e} z$ and $w \equiv_{f} z$.
By (MP8), $w' \equiv_{e + f f} z$.
It follows that $[w]_{e + f} \leq [w']_{i}$.
We have therefore proved that $[w]_{e + f} = [x]_{e} \vee [y]_{f}$.

(E6) holds.
We show that if $[x]_{e} \approx [y]_{f}$ then 
$s \bullet ([x]_{e} \vee [y]_{f}) = s \bullet [x]_{e} \vee s \bullet [y]_{f}$.
Let $s = s^{+}\hat{s}$.
Then $s \bullet ([x]_{e} \vee [y]_{f}) = [\hat{s} \cdot z]_{(s(e + f))^{+}}$.
Also, $s \bullet [x]_{e} = [\hat{s} \cdot x]_{(se)^{+}}$ and $s \bullet [y]_{f} = [\hat{s} \cdot y]_{(sf)^{+}}$.
From $z \equiv_{e} x$ and $z \equiv_{f} y$ and using (MPA6),
we have that $\hat{s} \cdot z \equiv_{\hat{s} \ast e} \hat{s} \cdot x$ and $\hat{s} \cdot z \equiv_{\hat{s} \ast f} \hat{s} \cdot y$.
Observe that $(se)^{+} \leq (\hat{s}e)^{+}$ and $(sf)^{+} \leq (\hat{s}f)^{+}$.
Thus by (MPA4), we have that 
$\hat{s} \cdot z \equiv_{(se)^{+}} \hat{s} \cdot x$ and $\hat{s} \cdot z \equiv_{(sf)^{+}} \hat{s} \cdot y$.
The result now follows, from the following calculation
$$(s(e+f))^{+} = (se \vee sf)^{+} = (se)^{+} + (sf)^{+}.$$

(E7) holds.
Let $s \vee t = (s \vee t)^{+}\hat{u}$ where $\hat{u} \in M$.
Thus $s \vee t = s^{+}\hat{u} \vee t^{+}\hat{u}$.
But $s = s^{+}\hat{u}$ and $t = t^{+}\hat{u}$.
The result now follows.

Let $X$ and $X'$ be $[B|M]$-sets.
Let $\alpha \colon X \rightarrow X'$ be a homomorphism of $[B|M]$-sets.
Put $Y = \bigcup_{e \in \mathsf{Proj}(S)} X/\equiv_{e}$
and
$Y' = \bigcup_{e \in \mathsf{Proj}(S)} X'/\equiv_{e}$.
Define $\theta \colon Y \rightarrow Y'$ by $\theta ([x]_{e}) = [\alpha (x)]_{e}$.
Then this is a homomorphism of supported actions by Theorem~\ref{them:mary}.
However, we also have to check that $\theta$ maps the smallest element of $Y$ to the smallest element of $Y'$
and that binary joins in $Y$ are preserved by $\theta$.
The proofs of these are both straightforward from the definition of $\alpha$.
We have therefore defined a functor from $[B|M]$-{\bf Set} to  $S$-{\bf FSup}.

To prove that we have an equivalence of categories, it is enough to prove that 
every supported action of a Boolean left restriction monoid is isomorphic to an induced supported action.
By Theorem~\ref{them:mary} there is an isomorphism,
but we need to check that minimum elements are mapped
and binary joins are preserved.
Both of these are easy.
\end{proof}



Let $S$ be a Boolean left restriction monoid.
Put $M = \mathsf{Tot}(S)$ and $B = \mathsf{Proj}(S)$.
Although not strictly necessary, we now sketch out the proof that the exponential object 
in the category $S$-{\bf FSup} gives rise to the exponential object in the category $[\mathsf{Proj}(S)|\mathsf{Tot}(S)]$-{\bf Set},
as described in \cite[Proposition 7.11]{G2024},
under the functor described in the theorem above.
Suppose that $(S,X,p)$ and $(S,Y,q)$ are factorizable supported actions.
We have that $(Y^{X})_{1} = \mbox{hom}(S \, \Box \, X,Y)$.
The set $S \,\Box \, X$ consists of those ordered pairs $(s,x)$ where $s^{+} = p(x)$.
Let $\theta \in \mbox{hom}(S \, \Box \, X,Y)$.
Define $\theta'$ to be the restriction of $\theta$ to the set $M \times X_{1}$.
Observe that $\theta' \colon M \times X_{1} \rightarrow Y_{1}$.
We may describe $\theta$ in terms of $\theta'$.
We calculate $\theta (s,x)$.
In $X$, we have that $x = p(x) \cdot x'$ where $x' \in X_{1}$.
In $S$, we may write $s = s^{+}\hat{s}$ where $\hat{s} \in M$.
It follows that $\theta (s,x) = s^{+} \theta (\hat{s},x')$ where $(\hat{s}, x') \in M \times X_{1}$.
Thus $\theta (s,x) = s^{+} \cdot \theta' (\hat{s}, x')$.
It is clear that $\theta'$ is a homomorphism of the $[B|M]$-sets $X_{1}$ and $Y_{1}$.
On the other hand, if $\phi' \colon M \times X_{1} \rightarrow Y_{1}$ is a homomorphism of $[B|M]$-sets,
define $\phi \colon S \, \Box \, X \rightarrow Y$
by $\phi (s,x) = s^{+} \cdot \phi' (\hat{s}, x')$. 
We need to show that $\phi$ is well-defined.
Suppose that $s = s^{+}\hat{s} = s^{+}\hat{t}$, where $\hat{s}, \hat{t} \in M$,
and $x = p(x) \cdot x' = p(x) \cdot x''$, where $x',x'' \in X_{1}$. 
Then $(\hat{s},x') \equiv_{s^{+}} (\hat{t}, x'')$.
But then we have that $\phi' (\hat{s},x') \equiv_{s^{+}} \phi' (\hat{t}, x'')$,
and the definition of $\equiv_{e}$ gives us the result.
In this way, we have established a bijection between the elements of $(Y^{X})_{1}$
and the elements of $M \times X_{1} \rightarrow Y_{1}$. 
This is the hard part of proving our result.

  


%%%%%%%%%%%%%%%%%%%%%%%%%%%%%%%%%%%%%%%%%%%%%%%%%%%%%%%%%%%%%%%%%%%%%%%%%%%%%%%%%%%%%%%%%%%%%%%%%%%%%%%%%%%%%%%%%%%%%%%%%%%%%%%%%%%%%%%%%
\section{Supported actions of \'etale Boolean restriction monoids}

Let $S$ be a Boolean left restriction monoid.
An element $a$ is called a {\em partial unit} if there is an element $b$ such that
$ab = a^{+}$ and $ba = b^{+}$.
The set of partial units of $S$ is denoted by $\mathsf{Inv}(S)$.
This subset forms a Boolean inverse monoid that contains all the projections of $S$.
A Boolean left restriction monoid is said to be  {\em \'etale} if every element of $S$
is a right-compatible join of partial units.
See \cite{L2025}, for a whole paper dedicated to the structure of \'etale Boolean left restriction monoids.

We shall compare supported actions of $S$ and supported actions of $\mathsf{Inv}(S)$.
Clearly, if there is a supported action of $S$ on $Y$ then there is a supported action of 
$\mathsf{Inv}(S)$ on $Y$.
We shall prove  the converse: that a supported action of $\mathsf{Inv}(S)$ on $Y$, induces a unique supported action of
$S$ on $Y$.

Let $S$ be a Boolean left restriction monoid and let $a \in S$.
Then, by assumption, there exist $a_{i} \in \mathsf{Inv}(S)$ such that
$a = \bigvee_{i=1}^{m} a_{i}$, where $a_{i} \sim_{r} a_{j}$.
It is therefore natural to define
$$a \cdot y = \bigvee_{i=1}^{m} a_{i} \cdot y.$$
The first step in showing that this is well-defined is to show that the join on the right hand side is actually defined in $Y$.

\begin{lemma}\label{lem:alice} Let $S$ be a Boolean left restriction monoid.
Suppose that $(\mathsf{Inv}(S),Y,p)$ is a supported action of a Boolean inverse monoid.
If $a_{i} \sim_{r} a_{j}$ and $y \in Y$ then $a_{i} \cdot y \approx a_{j} \cdot y$.
\end{lemma}
\begin{proof} We are given that $a_{i}^{+}a_{j} = a_{j}^{+}a_{i}$.
We shall calculate $p(a_{i} \cdot y) \cdot (a_{j} \cdot y)$.
Using the properties of the support function $p$, we have that
$$p(a_{i} \cdot y) \cdot (a_{j} \cdot y)
= 
(a_{i}p(y))^{+}a_{j} \cdot y.$$
Now observe that $a_{i}p(y) \leq a_{i}$ and so
$(a_{i}p(y))^{+} \leq a_{i}^{+}$.
It follows that
$(a_{i}p(y))^{+}a_{j} = (a_{i}p(y))^{+}a_{i}^{+}a_{j}$.
However, we know that $a_{i}^{+}a_{j} = a_{j}^{+}a_{i}$, 
and so
$(a_{i}p(y))^{+}a_{i}^{+}a_{j}
=
(a_{i}p(y))^{+}a_{j}^{+}a_{i}$.
But projections commute, and so
$(a_{i}p(y))^{+}a_{j}^{+}a_{i}
=
a_{j}^{+}(a_{i}p(y))^{+}a_{i}$
We now invoke axiom (LR6), to get
$a_{j}^{+}(a_{i}p(y))^{+}a_{i}
= 
a_{j}^{+}a_{i}p(y)$.
But $a_{j}^{+}a_{i}p(y) = a_{i}^{+}a_{j}p(y)$.
This proves that $p(a_{i} \cdot y) \cdot (a_{j} \cdot y) = a_{i}^{+}a_{j} \cdot y$.
The result now follows.
\end{proof}

We now show that our definition does not depend on our choice of elements from $\mathsf{Inv}(S)$.

\begin{lemma} With the above notation, if $a = \bigvee_{i=1}^{m} a_{i}$ and $a = \bigvee_{j = 1}^{n} b_{j}$,
where $a_{i}, b_{j} \in \mathsf{Inv}(S)$,
then $\bigvee_{i=1}^{m} a_{i} \cdot y = \bigvee_{j=1}^{n} b_{j} \cdot y$.
\end{lemma}
\begin{proof} We have that $a_{k} = \bigvee_{j=1}^{n} a_{k}^{+}b_{j}$ for each $1 \leq k \leq m$.
Observe that $(a_{i}^{+}b_{j}) \cdot y \leq b_{j} \cdot y$.
We deduce that $\bigvee_{i=1}^{m} a_{i} \cdot y \leq \bigvee_{j=1}^{n} b_{j} \cdot y$.
Symmetry delivers the proof of the other direction.
\end{proof}

Our definition is therefore well-defined.
It is easy to check that axioms (E1) and (E2) hold.
All the axioms (E3)--(E6) are easy to show.
We prove (E7).

\begin{lemma} With the above definitions, if $a \sim_{r} b$ then $(a \vee b) \cdot y = a \cdot y \vee b \cdot y$.
\end{lemma}
\begin{proof} Suppose that $a = \bigvee_{i=1}^{m} a_{i}$ and $b = \bigvee_{j=1}^{n} b_{j}$ where
$a_{i}, b_{j} \in \mathsf{Inv}(S)$.
The first step is to prove  that $a_{i} \sim_{r} b_{j}$.
We are given that $b^{+}a = a^{+}b$.
But $a_{i} \leq a$.
Thus $b^{+}a_{i} \leq a^{+}b$.
But $b_{j} \leq b$.
Thus $b_{j}^{+}a_{i} \leq a^{+}b$.
Mutiply both sides by $a_{i}^{+}$ and use the fact that projections commute.
Thus $b_{j}^{+}a_{i} \leq a_{i}^{+}a^{+}b$.
But $a_{i}^{+} \leq a^{+}$.
Thus $b_{j}^{+}a_{i} \leq a_{i}^{+}b$.	 
Multiply both sides by $b_{j}^{+}$
and use the fact that projections commute.
Thus $b_{j}^{+}a_{i} \leq a_{i}^{+}b_{j}^{+}b$.	 
However, $b_{j} = b_{j}^{+}b$.
We have therefore show that 
$b_{j}^{+}a_{i} \leq a_{i}^{+}b_{j}$.	
The result now follows by symmetry. 
By Lemma~\ref{lem:alice}, this implies that $a_{i} \cdot y \approx b_{j} \cdot y$.
This means that $a_{i} \cdot y \vee b_{j} \cdot y$ is defined.
It follows that $a \cdot y \vee b \cdot y$ is defined.
The element $(a \vee b) \cdot y$ is defined by assumption.
It remains to prove equality.
By definition, $(a \vee b) \cdot y = \left( \bigvee  a_{i} \cdot y \right) \vee \left( \bigvee b_{j} \cdot y \right) = \bigvee a_{i} \cdot y \vee b_{j} \cdot y$.
The result will therefore be proved if we show that $(a_{i} \vee b_{j}) \cdot y = a_{i} \cdot y \vee b_{j} \cdot y$.
We shall prove this result by way of the following.
Let $s,t \in \mathsf{Inv}(S)$.
Supppose that $s \sim_{r} t$.
We prove that $(s \vee t) \cdot y = s \cdot y \vee t \cdot y$.
It is immediate that $s \cdot y \vee t \cdot y \leq (s \vee t) \cdot y$.
The result now follows when we observe that $y = p(y) \cdot y$.
\end{proof}


The proof of the following is now almost immediate by what we have proved in this section.

\begin{theorem} Let $S$ be an \'etale Boolean left restriction monoid.
Then the category of supported actions of $\mathsf{Inv}(S)$ is isomorphic
with the category of supported actions of $S$.
\end{theorem}










%%%%%%%%%%%%%%%%%%%%%%%%%%%%%%%%%%%%%%%%%%%%%%%%%%%%%%%%%%%%%%%%%%%%%%%%%%%%%%%%%%%%%%%%%%%%%%%%%%%%%%%%%%%%%%%%%%%%%%%%%%%%%%%%%%%%%%%%%%%%%%%%%%%%%%%%%%%%%%%%%%%%%%%%%%%%%%%%%%%%%
\begin{thebibliography}{99}


\bibitem{BW} M. Barr, Ch. Wells, {\em Category theory for computing science}, Prentice Hall, 1990.

%\bibitem{B1991} G. M. Bergman, Actions of Boolean rings on sets,
%{\em Algebra Universalis} {\bf 28} (1991), 153--187.


\bibitem{EF} D. Easdown, D. G. FitzGerald, Presentations of factorizable inverse monoids, 
{\em Acta Sci. Math. (Szeged)} {\bf 71} (2005), 509--520.

\bibitem{F2010} J. Funk, B. Steinberg, The universal covering of an inverse semigroup,
{\em Appl. Categ. Structures} {\bf 18} (2010), 135--163.

\bibitem{G2024} R. Garner, Cartesian closed varieties I: the classification theorem, arXiv:2302.04402.

\bibitem{G2025} R. Garner, Cartesian closed varieties II: links to operator algebra, arXiv:2302.04403.

\bibitem{JS} M. Jackson, T. Stokes, Modal restriction semigroups: towards an algebra of functions,
{\em Int. J. Algebra Comput.} {\bf 21} (2011), 1053--1095.

%\bibitem{J1990} P. T. Johnstone, Collapsed toposes and Cartesian closed varieties, 
%{\em J. Algebra} {\bf 129} (1990), 446--480.
tube
\bibitem{L1988} J. Lambek, P. J. Scott, {\em Introduction to higher order categorical logic}, CUP, Cambridge, 1986.

\bibitem{L2025} M. V. Lawson, The structure of \'etale Boolean right restriction monoids, arXiv:2404.08606.

\bibitem{LR2021} M. V. Lawson, P. Resende, Morita equivalence of pseudogroups, {\em J. Algebra} {\bf 586} 2021, 718--755.

\bibitem{S2011} B. Steinberg, Strong Morita equivalence of inverse semigroups,
{\em Houston J. Math.} {\bf 37} (2011), 895--927.




\end{thebibliography}
\end{document}





%%%%%%%%%%%%%%%%%%%%%%%%%%%%%%%%%%%%%%%%%%%%%%%%%%%%%%%%%%%%%%%%%%%%%%%%%%%%%%%%%%%%%%%%%%%%%%%%%%%%%%%%%%%%%%%%%%%%%%%%%%%%%%%%%%%%%%%%%%%%%%%%%%%%%


%%%%%%%%%%%%%%%%%%%%%%%%%%%%%%%%%%%%%%%%%%%%%%%%%%%%%%%%%%%%%%%%%%%%%%%%%%%%%%%%%%%%%%%%%%%%%%%%%%%%%%%%%%%%%%%%%%%%%%%%%%%%%%%%%%%%%%%%%%%%%%%%%%%%%

The following result uses \cite{B1991} combined with the remark in Section~3.

\begin{lemma} Let $(S,X,p)$ be a supported action of a Boolean left restriction monoid where $X_{1}$ is non-empty.
Then $X$ is a sheaf over a Boolean algebra.
\end{lemma}
\begin{proof} Let $e$ be any projection and let $f \leq e$.
Then $e = f + e\bar{f}$.
There are two maps $\phi^{e}_{f} \colon X_{e} \rightarrow X_{f}$ and $\phi^{e}_{e\bar{f}} \colon X_{e} \rightarrow X_{e\bar{f}}$.
If $x \in X_{e}$ then $x = fx \vee e\bar{f}x$
since $fx \approx e\bar{f}x$.
This enable us to define a bijection from $X_{e}$ to $X_{f} \times X_{e\bar{f}}$.
We have therefore verified \cite[Proposition 1]{B1991}.
\end{proof}















The result above suggests that Boolean left restriction monoids can be constructed out of Boolean algebras and (smaller) monoids.
How this is to be done was discovered by Garner.
To describe his results, we shall need some definitions.

Let $B$ be a Boolean algebra.
A {\em $B$-set}, a notion first defined in \cite{B1991}, is a set $X$ for which there is ternary operation $B \times X \times X \rightarrow X$ denoted by $(e,x,y) \mapsto e(x,y)$
satisfying the following axioms:
\begin{itemize}
\item[{\rm (BS1).}] $e(x,x) = x$.
\item[{\rm (BS2).}] $e(e(x,y),z) = e(x,z)$.
\item[{\rm (BS3).}] $e(x,e(y,z)) = e(x,z)$.
\item[{\rm (BS4).}] $1(x,y) = x$.
\item[{\rm (BS5).}] $\bar{e}(x,y) = e(y,x)$.
\item[{\rm (BS6).}] $ef(x,y) = e(f(x,y),y)$. 
\end{itemize}

\begin{remark}{\em Observe that $B$ does not really act on the set $X$.
The goal of this paper can be viewed as obtaining an honest-to-goodness action.}
\end{remark}


If $X$ is a $B$-set, then for each $e \in B$ we may define an equivalence relation $\equiv_{e}$ on $X$
by $x \equiv_{e} y$ iff $e(x,y) = y$.
Such equivalence relations are the subject of \cite[Proposition 3.2]{G2024}.
If $X$ and $X'$ are $B$-sets then a homomorphism is a function $\theta \colon X \rightarrow X'$ such that
$\theta((e(x,y)) = e(\theta(x),\theta(y))$ for all $e \in B$.

Let $M$ be a monoid.
By an {\em $M$-set} we shall always mean a left action of the monoid $M$ on a set.

A {\em matched pair} $[B|M]$, defined first as \cite[Definition 7.1]{G2024}, consists of a Boolean algebra $B$ and a monoid $M$ satisfying the following axioms:
\begin{itemize}
\item[{\rm (MP1).}] $M$ is a $B$-set.
\item[{\rm (MP2).}] $B$ is an $M$-set with the action denoted by $(a,e) \mapsto a \ast e$.
\item[{\rm (MP3).}] $M$ acts on $B$ via Boolean algebra homomorphisms.
\item[{\rm (MP4).}] $e(m,n)p = e(mp,np)$.
\item[{\rm (MP5).}] $m(e(n,p)) = (m \ast e)(mn,mp)$.
\item[{\rm (MP6).}] $e(m,n) \ast f = e(m \ast f, n \ast f)$.
\end{itemize}
If $[B|M]$ and $[B'|M']$ are matched pairs then a homomorphism from $[B|M]$ to $[B'|M']$ is a pair
$(\phi, \theta)$ where $\theta \colon M \rightarrow M'$ is a monoid homomorphism,
$\phi \colon B \rightarrow B'$ is a homomorphism of Boolean algebras
and $\phi (m \ast e) = \theta(m) \ast \phi (e)$ and $\theta (e(m,n)) = \phi (e) (\theta (m), \theta (n))$.

Each Boolean left restriction monoid $S$ gives rise to a matched pair as follows.
Put $B = \mathsf{Proj}(S)$ and $M = \mathsf{Tot}(S)$.
Define an action of $M$ on $B$ by 
$$a \ast e = (ae)^{+}$$
and define a $B$-structure on $M$ by $(e,a,b) \mapsto ea \vee \bar{e}b$.
It is now easy to check that $[B|M]$ is a matched pair.
In \cite[Proposition 3.3, Proposition 3.4, Theorem 3.5]{G2025}, Garner proves that matched pairs determine the structure of Boolean left restriction monoids.
We prove the following result which is implicit in \cite{G2025}.


\begin{proposition}\label{prop:one}
Let $S$ and $S'$ be Boolean left restriction monoids,
with associated matched pairs $[B|M]$ and $[B'|M']$, respectively, so that
$S = S[B|M]$ and $S' = S'[B'|M']$.
Then if $[B|M]$ and $[B'|M']$ are isomorphic as matched pairs
the semigroups $S$ and $S'$ are isomorphic as Boolean left restriction monoids.
\end{proposition}
\begin{proof} Let $(\phi, \theta)$ be the isomorphism from the matched pair $[B|M]$ to the matched pair $[B'|M']$.
If $s \in S$ then as can write $s = s^{+}\hat{s}$, where $\hat{s}$ is any total element of $S$.
Define $\Theta \colon S \rightarrow S'$ by $\theta (s) = \phi (s^{+})\theta(\hat{s})$. 
This is well-defined, since if $s = s^{+}\hat{t}$ then $\hat{s} \equiv_{s^{+}} \hat{t}$ and so $\theta (\hat{s}) \equiv_{\phi (s^{+})} \theta (\hat{t})$.
The function $\Theta$ is injective, since if $\Theta (s) = \Theta (t)$,
then $\phi (s^{+})f(\hat{s}) = \phi (t^{+}) f(\hat{t})$.
We now use axiom (LR5) to deduce that $\phi (s^{+}) = \phi (t^{+})$ which gives $s^{+} = t^{+}$.
As a consequence, $\theta (\hat{s}) \equiv_{\phi (s^{+})} \theta (\hat{t})$ from which we deduce that $\hat{s} \equiv_{s^{+}} \hat{t}$.
Thus $s = t$.
The fact that $\Theta$ is surjective is immediate.
That $\Theta$ is a homomorphism follows from the fact that if $s \leq \hat{s}$ and $t \leq \hat{t}$ then $st \leq \hat{s}\hat{t}$,
where $\hat{s}\hat{t}$ is a total map.
Thus we have proved that $\Theta$ is an isomorphism of semigroups.
It is an isomorphism of left restriction monoids because $\Theta (s^{+}) = \Theta (s)^{+}$.
We claim that $\Theta$ is an isomorphism of Boolean left restriction monoids.
To see why,
suppose that $s \vee t$ is defined in $S$.
This element is beneath a total element, which we shall denote by $\hat{v}$.
We may therefore write $s \vee t = (s \vee t)^{+} \hat{v}$.
using Lemma~\ref{lem:induction} and distibutivity, we have that
$s \vee t = s^{+}\hat{t} \vee t^{+}\hat{t}$.
Observe that $s, s^{+}\hat{v} \leq \hat{v}$.
Thus $s \sim_{r} s^{+}\hat{v}$.
By axiom (LR5), we can see that $s^{+} = (s^{+}\hat{v})^{+}$.
Thus $s = s^{+}\hat{t}$ and, similarly, $t = t^{+}\hat{v}$.
It follows that $\Theta (s) = \phi (s^{+})\theta (\hat{v})$, $\Theta (t) = \phi (t^{+}) \theta (\hat{v})$
and $\Theta (s \vee t) = \phi ((s \vee t)^{+}) \theta (\hat{v})$.
But $\phi$ is a homorphism of Boolean algebras.
The claim is therefore proved.
\end{proof}

We shall now explain what we prove in this paper.
To do so will require another definition.
Let $[B|M]$ be a matched pair.
A {\em $[B|M]$-set} is a set $X$ that satisfies the following axioms:
\begin{itemize}
\item[{\rm (BMS1).}] $X$ is a $B$-set.
\item[{\rm (BMS2).}] $X$ is an $M$-set where the action is denoted by $(m,x) \mapsto m \cdot x$.
\item[{\rm (BMS3).}] $e(m,n) \cdot x = e(m \cdot x, n \cdot x)$.
\item[{\rm (BMS4).}] $m \cdot e(x,y) = (m \ast e)(m \cdot x, m \cdot y)$.
\end{itemize}
A homomorphism of $[B|M]$-sets is a function that is both a $B$-set homomorphism and an $M$-set homomorphism.
We may therefore form the category of $[B|M]$-sets which we denote by {\bf $[B|M]$-sets}. 
The goal of this paper can now be stated: {\em we shall prove that the category {\bf $[B|M]$-sets} is isomorphic to a category of special actions of the monoid $S[B|M]$.}
Just what those actions are is defined in the next section.






\begin{lemma} Let $(S,X)$ and $(S,X')$ be \'etale actions of the left restriction monoid $S$.
Let $\theta \colon X \rightarrow X'$ be a homomorphism.
Then $\theta$ is order preserving.
Furthermore, 
if $\theta$ is a bijection then $\theta$ is an order isomorphism.
\end{lemma}
\begin{proof}
Let $y_{1}, y_{2} \in Y$ such that $y_{1} \leq y_{2}$.
Then $y_{1} = e \cdot y_{2}$ where $e$ is a projection.
Then $\theta (y_{1}) = e \cdot \theta (y_{2})$.
It follows that $\theta (y_{1}) \leq \theta (y_{2})$.
Thus $\theta$ is order preserving.

Now suppose that $\theta$ is invertible.
Let $y' \in Y'$ and $a \in S$.
Then $\theta (y) = y'$ for a unique element $y \in Y$.
Now, $\theta (s \cdot y) = s \cdot \theta (y) = s \cdot y'$.
It follows that $\theta^{-1}(s \cdot y') = s \cdot y = s \cdot \theta^{-1}(y')$.
Clearly, $p' = p \theta^{-1}$.
Thus $\theta^{-1}$ is a homomorphism and so order preserving.
\end{proof}









